La medición de cobertura mencionada en \cref{SEC:TEST-COVERAGE} se
desglosa por aplicación Django en la \cref{table:cdetail-coverage}. % TODO cuadro en vez de tabla

\begin{table}[Cobertura del código por aplicación]%
            {table:cdetail-coverage}%
            {Cobertura del código de aplicación por \textit{app} Django,
             excluyendo tests y migraciones.}
  \small
  \renewcommand{\arraystretch}{1.3} % TODO fix percentages
  \begin{tabular}{l r r r}
    \hline
    \textbf{App} & \textbf{Sentencias} & \textbf{No cubiertas} & \textbf{Cobertura} \\
    \hline
    app1        & lines & uncovered & coverage\,\% \\
    \hline
    \textbf{Total}  & \textbf{1} & \textbf{1} & \textbf{88\,\%} \\
    \hline
  \end{tabular}
\end{table}

La columna ``Cobertura'' refleja el porcentaje de sentencias del código
de aplicación efectivamente ejecutadas por las pruebas. La cifra agregada del
informe que incluye estos archivos asciende al 88\,\%
(\cref{SEC:TEST-COVERAGE}).
